\IfFileExists{stacks-project.cls}{%
\documentclass{stacks-project}
}{%
\documentclass{amsart}
}

% For dealing with references we use the comment environment
\usepackage{verbatim}
\newenvironment{reference}{\comment}{\endcomment}
%\newenvironment{reference}{}{}
\newenvironment{slogan}{\comment}{\endcomment}
\newenvironment{history}{\comment}{\endcomment}

% For commutative diagrams we use Xy-pic
\usepackage[all]{xy}

% We use 2cell for 2-commutative diagrams.
\xyoption{2cell}
\UseAllTwocells

% We use multicol for the list of chapters between chapters
\usepackage{multicol}

% This is generally recommended for better output
\usepackage{lmodern}
\usepackage[T1]{fontenc}

% For cross-file-references
\usepackage{xr-hyper}

% Package for hypertext links:
\usepackage{hyperref}

% For any local file, say "hello.tex" you want to link to please
% use \externaldocument[hello-]{hello}
\externaldocument[introduction-]{introduction}
\externaldocument[conventions-]{conventions}
\externaldocument[sets-]{sets}
\externaldocument[categories-]{categories}
\externaldocument[topology-]{topology}
\externaldocument[sheaves-]{sheaves}
\externaldocument[sites-]{sites}
\externaldocument[stacks-]{stacks}
\externaldocument[fields-]{fields}
\externaldocument[algebra-]{algebra}
\externaldocument[brauer-]{brauer}
\externaldocument[homology-]{homology}
\externaldocument[derived-]{derived}
\externaldocument[simplicial-]{simplicial}
\externaldocument[more-algebra-]{more-algebra}
\externaldocument[smoothing-]{smoothing}
\externaldocument[modules-]{modules}
\externaldocument[sites-modules-]{sites-modules}
\externaldocument[injectives-]{injectives}
\externaldocument[cohomology-]{cohomology}
\externaldocument[sites-cohomology-]{sites-cohomology}
\externaldocument[dga-]{dga}
\externaldocument[dpa-]{dpa}
\externaldocument[sdga-]{sdga}
\externaldocument[hypercovering-]{hypercovering}
\externaldocument[schemes-]{schemes}
\externaldocument[constructions-]{constructions}
\externaldocument[properties-]{properties}
\externaldocument[morphisms-]{morphisms}
\externaldocument[coherent-]{coherent}
\externaldocument[divisors-]{divisors}
\externaldocument[limits-]{limits}
\externaldocument[varieties-]{varieties}
\externaldocument[topologies-]{topologies}
\externaldocument[descent-]{descent}
\externaldocument[perfect-]{perfect}
\externaldocument[more-morphisms-]{more-morphisms}
\externaldocument[flat-]{flat}
\externaldocument[groupoids-]{groupoids}
\externaldocument[more-groupoids-]{more-groupoids}
\externaldocument[etale-]{etale}
\externaldocument[chow-]{chow}
\externaldocument[intersection-]{intersection}
\externaldocument[pic-]{pic}
\externaldocument[weil-]{weil}
\externaldocument[adequate-]{adequate}
\externaldocument[dualizing-]{dualizing}
\externaldocument[duality-]{duality}
\externaldocument[discriminant-]{discriminant}
\externaldocument[derham-]{derham}
\externaldocument[local-cohomology-]{local-cohomology}
\externaldocument[algebraization-]{algebraization}
\externaldocument[curves-]{curves}
\externaldocument[resolve-]{resolve}
\externaldocument[models-]{models}
\externaldocument[functors-]{functors}
\externaldocument[equiv-]{equiv}
\externaldocument[pione-]{pione}
\externaldocument[etale-cohomology-]{etale-cohomology}
\externaldocument[proetale-]{proetale}
\externaldocument[relative-cycles-]{relative-cycles}
\externaldocument[more-etale-]{more-etale}
\externaldocument[trace-]{trace}
\externaldocument[crystalline-]{crystalline}
\externaldocument[spaces-]{spaces}
\externaldocument[spaces-properties-]{spaces-properties}
\externaldocument[spaces-morphisms-]{spaces-morphisms}
\externaldocument[decent-spaces-]{decent-spaces}
\externaldocument[spaces-cohomology-]{spaces-cohomology}
\externaldocument[spaces-limits-]{spaces-limits}
\externaldocument[spaces-divisors-]{spaces-divisors}
\externaldocument[spaces-over-fields-]{spaces-over-fields}
\externaldocument[spaces-topologies-]{spaces-topologies}
\externaldocument[spaces-descent-]{spaces-descent}
\externaldocument[spaces-perfect-]{spaces-perfect}
\externaldocument[spaces-more-morphisms-]{spaces-more-morphisms}
\externaldocument[spaces-flat-]{spaces-flat}
\externaldocument[spaces-groupoids-]{spaces-groupoids}
\externaldocument[spaces-more-groupoids-]{spaces-more-groupoids}
\externaldocument[bootstrap-]{bootstrap}
\externaldocument[spaces-pushouts-]{spaces-pushouts}
\externaldocument[spaces-chow-]{spaces-chow}
\externaldocument[groupoids-quotients-]{groupoids-quotients}
\externaldocument[spaces-more-cohomology-]{spaces-more-cohomology}
\externaldocument[spaces-simplicial-]{spaces-simplicial}
\externaldocument[spaces-duality-]{spaces-duality}
\externaldocument[formal-spaces-]{formal-spaces}
\externaldocument[restricted-]{restricted}
\externaldocument[spaces-resolve-]{spaces-resolve}
\externaldocument[formal-defos-]{formal-defos}
\externaldocument[defos-]{defos}
\externaldocument[cotangent-]{cotangent}
\externaldocument[examples-defos-]{examples-defos}
\externaldocument[algebraic-]{algebraic}
\externaldocument[examples-stacks-]{examples-stacks}
\externaldocument[stacks-sheaves-]{stacks-sheaves}
\externaldocument[criteria-]{criteria}
\externaldocument[artin-]{artin}
\externaldocument[quot-]{quot}
\externaldocument[stacks-properties-]{stacks-properties}
\externaldocument[stacks-morphisms-]{stacks-morphisms}
\externaldocument[stacks-limits-]{stacks-limits}
\externaldocument[stacks-cohomology-]{stacks-cohomology}
\externaldocument[stacks-perfect-]{stacks-perfect}
\externaldocument[stacks-introduction-]{stacks-introduction}
\externaldocument[stacks-more-morphisms-]{stacks-more-morphisms}
\externaldocument[stacks-geometry-]{stacks-geometry}
\externaldocument[moduli-]{moduli}
\externaldocument[moduli-curves-]{moduli-curves}
\externaldocument[examples-]{examples}
\externaldocument[exercises-]{exercises}
\externaldocument[guide-]{guide}
\externaldocument[desirables-]{desirables}
\externaldocument[coding-]{coding}
\externaldocument[obsolete-]{obsolete}
\externaldocument[fdl-]{fdl}
\externaldocument[index-]{index}

% Theorem environments.
%
\theoremstyle{plain}
\newtheorem{theorem}[subsection]{Theorem}
\newtheorem{proposition}[subsection]{Proposition}
\newtheorem{lemma}[subsection]{Lemma}

\theoremstyle{definition}
\newtheorem{definition}[subsection]{Definition}
\newtheorem{example}[subsection]{Example}
\newtheorem{exercise}[subsection]{Exercise}
\newtheorem{situation}[subsection]{Situation}

\theoremstyle{remark}
\newtheorem{remark}[subsection]{Remark}
\newtheorem{remarks}[subsection]{Remarks}

\numberwithin{equation}{subsection}

% Macros
%
\def\lim{\mathop{\mathrm{lim}}\nolimits}
\def\colim{\mathop{\mathrm{colim}}\nolimits}
\def\Spec{\mathop{\mathrm{Spec}}}
\def\Hom{\mathop{\mathrm{Hom}}\nolimits}
\def\Ext{\mathop{\mathrm{Ext}}\nolimits}
\def\SheafHom{\mathop{\mathcal{H}\!\mathit{om}}\nolimits}
\def\SheafExt{\mathop{\mathcal{E}\!\mathit{xt}}\nolimits}
\def\Sch{\mathit{Sch}}
\def\Mor{\mathop{\mathrm{Mor}}\nolimits}
\def\Ob{\mathop{\mathrm{Ob}}\nolimits}
\def\Sh{\mathop{\mathit{Sh}}\nolimits}
\def\NL{\mathop{N\!L}\nolimits}
\def\CH{\mathop{\mathrm{CH}}\nolimits}
\def\proetale{{pro\text{-}\acute{e}tale}}
\def\etale{{\acute{e}tale}}
\def\QCoh{\mathit{QCoh}}
\def\Ker{\mathop{\mathrm{Ker}}}
\def\Im{\mathop{\mathrm{Im}}}
\def\Coker{\mathop{\mathrm{Coker}}}
\def\Coim{\mathop{\mathrm{Coim}}}

% Boxtimes
%
\DeclareMathSymbol{\boxtimes}{\mathbin}{AMSa}{"02}

%
% Macros for moduli stacks/spaces
%
\def\QCohstack{\mathcal{QC}\!\mathit{oh}}
\def\Cohstack{\mathcal{C}\!\mathit{oh}}
\def\Spacesstack{\mathcal{S}\!\mathit{paces}}
\def\Quotfunctor{\mathrm{Quot}}
\def\Hilbfunctor{\mathrm{Hilb}}
\def\Curvesstack{\mathcal{C}\!\mathit{urves}}
\def\Polarizedstack{\mathcal{P}\!\mathit{olarized}}
\def\Complexesstack{\mathcal{C}\!\mathit{omplexes}}
% \Pic is the operator that assigns to X its picard group, usage \Pic(X)
% \Picardstack_{X/B} denotes the Picard stack of X over B
% \Picardfunctor_{X/B} denotes the Picard functor of X over B
\def\Pic{\mathop{\mathrm{Pic}}\nolimits}
\def\Picardstack{\mathcal{P}\!\mathit{ic}}
\def\Picardfunctor{\mathrm{Pic}}
\def\Deformationcategory{\mathcal{D}\!\mathit{ef}}
\usepackage{cleveref}
\crefname{originalreference}{original reference}{original references}
\Crefname{originalreference}{Original reference}{Original references}
\makeatletter
\newlabel{algebra-lemma-ring-colimit-fp}{{00QN}{source}{}{sourcepointer-1}{}}
\newlabel{algebra-lemma-ring-colimit-fp@cref}{{[originalreference][0][]00QN}{[1][0][]source}{}{}{}}
\makeatother
\crefname{originalreference}{original reference}{original references}
\Crefname{originalreference}{Original reference}{Original references}
\makeatletter
\newlabel{algebraic-lemma-characterize-representable-by-space}{{02ZX}{source}{}{sourcepointer-2}{}}
\newlabel{algebraic-lemma-characterize-representable-by-space@cref}{{[originalreference][0][]02ZX}{[1][0][]source}{}{}{}}
\makeatother
\crefname{originalreference}{original reference}{original references}
\Crefname{originalreference}{Original reference}{Original references}
\makeatletter
\newlabel{bootstrap-lemma-glueing-sheaves}{{07WE}{source}{}{sourcepointer-3}{}}
\newlabel{bootstrap-lemma-glueing-sheaves@cref}{{[originalreference][0][]07WE}{[1][0][]source}{}{}{}}
\makeatother
\crefname{originalreference}{original reference}{original references}
\Crefname{originalreference}{Original reference}{Original references}
\makeatletter
\newlabel{sets-lemma-bound-size-fibre-product}{{04T6}{source}{}{sourcepointer-4}{}}
\newlabel{sets-lemma-bound-size-fibre-product@cref}{{[originalreference][0][]04T6}{[1][0][]source}{}{}{}}
\makeatother
\crefname{originalreference}{original reference}{original references}
\Crefname{originalreference}{Original reference}{Original references}
\makeatletter
\newlabel{sets-lemma-what-is-in-it}{{000R}{source}{}{sourcepointer-5}{}}
\newlabel{sets-lemma-what-is-in-it@cref}{{[originalreference][0][]000R}{[1][0][]source}{}{}{}}
\makeatother
\crefname{originalreference}{original reference}{original references}
\Crefname{originalreference}{Original reference}{Original references}
\makeatletter
\newlabel{spaces-lemma-coproduct-algebraic-spaces}{{02WQ}{source}{}{sourcepointer-6}{}}
\newlabel{spaces-lemma-coproduct-algebraic-spaces@cref}{{[originalreference][0][]02WQ}{[1][0][]source}{}{}{}}
\makeatother
\crefname{originalreference}{original reference}{original references}
\Crefname{originalreference}{Original reference}{Original references}
\makeatletter
\newlabel{spaces-lemma-glueing-algebraic-spaces}{{02WR}{source}{}{sourcepointer-7}{}}
\newlabel{spaces-lemma-glueing-algebraic-spaces@cref}{{[originalreference][0][]02WR}{[1][0][]source}{}{}{}}
\makeatother
\crefname{originalreference}{original reference}{original references}
\Crefname{originalreference}{Original reference}{Original references}
\makeatletter
\newlabel{stacks-lemma-2-product-stacks-in-groupoids}{{02ZL}{source}{}{sourcepointer-8}{}}
\newlabel{stacks-lemma-2-product-stacks-in-groupoids@cref}{{[originalreference][0][]02ZL}{[1][0][]source}{}{}{}}
\makeatother
\title{A limit-preserving base with locally presentable fibres is representable}
\author{The Stacks Project Authors}
\date{stacks-project:criteria.tex:local-snapshot-20261012; GNU FDL 1.2 or later}
\begin{document}
\maketitle
\noindent\textit{Editorial note.} The footnote '(details omitted)' covers the set-theoretic size check of the glueing lemma; retained verbatim and flagged for the reviewer.\par\medskip
\noindent\textit{Editorial note.} Steps the author leaves to the reader are flagged here for the reviewer: source line(s) 824 carry an "omitted"/"follows easily"/"immediate from the definitions" phrase. Those sentences are retained verbatim; the displayed argument around them stays inside the counted span.\par\medskip
\noindent\textit{Allowed standard prerequisites (editorial).}\begin{itemize}
\item Affines are limits of finite presentation. algebra-lemma-ring-colimit-fp (cross-chapter, named). Writes V = lim V\_i in Case I.
\item Stackiness of fibre products. stacks-lemma-2-product-stacks-in-groupoids (same chapter, named). Gives Zariski stackiness of the fibre product in Case II.
\item Glueing sheaves and spaces. bootstrap-lemma-glueing-sheaves, spaces-lemma-glueing-algebraic-spaces (cross-chapter, named). Reassembles Z from its local pieces.
\end{itemize}\medskip
\ifdefined\hypertarget\hypertarget{sourcepointer-1}{}\fi
\noindent\textit{Source reference.} Original reference 00QN: focused-seeds/geometry/S9-crit-check-representable-limit-preserving/ref-tags.txt, source lines 5--5. Navigation pointer only: the counted proof cites this statement by \textbackslash{}ref as a standing input; the deduction the counted proof has to make is carried out inside the counted span, and no part of that deduction is delegated to this pointer.\par\medskip
\ifdefined\hypertarget\hypertarget{sourcepointer-2}{}\fi
\noindent\textit{Source reference.} Original reference 02ZX: focused-seeds/geometry/S9-crit-check-representable-limit-preserving/ref-tags.txt, source lines 6--6. Navigation pointer only: the counted proof cites this statement by \textbackslash{}ref as a standing input; the deduction the counted proof has to make is carried out inside the counted span, and no part of that deduction is delegated to this pointer.\par\medskip
\ifdefined\hypertarget\hypertarget{sourcepointer-3}{}\fi
\noindent\textit{Source reference.} Original reference 07WE: focused-seeds/geometry/S9-crit-check-representable-limit-preserving/ref-tags.txt, source lines 7--7. Navigation pointer only: the counted proof cites this statement by \textbackslash{}ref as a standing input; the deduction the counted proof has to make is carried out inside the counted span, and no part of that deduction is delegated to this pointer.\par\medskip
\ifdefined\hypertarget\hypertarget{sourcepointer-4}{}\fi
\noindent\textit{Source reference.} Original reference 04T6: focused-seeds/geometry/S9-crit-check-representable-limit-preserving/ref-tags.txt, source lines 8--8. Navigation pointer only: the counted proof cites this statement by \textbackslash{}ref as a standing input; the deduction the counted proof has to make is carried out inside the counted span, and no part of that deduction is delegated to this pointer.\par\medskip
\ifdefined\hypertarget\hypertarget{sourcepointer-5}{}\fi
\noindent\textit{Source reference.} Original reference 000R: focused-seeds/geometry/S9-crit-check-representable-limit-preserving/ref-tags.txt, source lines 9--9. Navigation pointer only: the counted proof cites this statement by \textbackslash{}ref as a standing input; the deduction the counted proof has to make is carried out inside the counted span, and no part of that deduction is delegated to this pointer.\par\medskip
\ifdefined\hypertarget\hypertarget{sourcepointer-6}{}\fi
\noindent\textit{Source reference.} Original reference 02WQ: focused-seeds/geometry/S9-crit-check-representable-limit-preserving/ref-tags.txt, source lines 10--10. Navigation pointer only: the counted proof cites this statement by \textbackslash{}ref as a standing input; the deduction the counted proof has to make is carried out inside the counted span, and no part of that deduction is delegated to this pointer.\par\medskip
\ifdefined\hypertarget\hypertarget{sourcepointer-7}{}\fi
\noindent\textit{Source reference.} Original reference 02WR: focused-seeds/geometry/S9-crit-check-representable-limit-preserving/ref-tags.txt, source lines 11--11. Navigation pointer only: the counted proof cites this statement by \textbackslash{}ref as a standing input; the deduction the counted proof has to make is carried out inside the counted span, and no part of that deduction is delegated to this pointer.\par\medskip
\ifdefined\hypertarget\hypertarget{sourcepointer-8}{}\fi
\noindent\textit{Source reference.} Original reference 02ZL: focused-seeds/geometry/S9-crit-check-representable-limit-preserving/ref-tags.txt, source lines 12--12. Navigation pointer only: the counted proof cites this statement by \textbackslash{}ref as a standing input; the deduction the counted proof has to make is carried out inside the counted span, and no part of that deduction is delegated to this pointer.\par\medskip
\begin{lemma}
\label{lemma-check-representable-limit-preserving}
Let $S$ be a scheme.
Let $\kappa = \text{size}(T)$ for some $T \in \Ob((\Sch/S)_{fppf})$.
Let $f : \mathcal{X} \to \mathcal{Y}$ be a $1$-morphism
of categories fibred in groupoids over $(\Sch/S)_{fppf}$
such that
\begin{enumerate}
\item $\mathcal{Y} \to (\Sch/S)_{fppf}$ is limit preserving on objects,
\item for an affine scheme $V$ locally of finite presentation over $S$ and
$y \in \Ob(\mathcal{Y}_V)$ the fibre product
$(\Sch/V)_{fppf} \times_{y, \mathcal{Y}} \mathcal{X}$ is representable
by an algebraic space of size $\leq \kappa$\footnote{The condition on
size can be dropped by those ignoring set theoretic issues.},
\item $\mathcal{X}$ and $\mathcal{Y}$ are stacks for the Zariski topology.
\end{enumerate}
Then $f$ is representable by algebraic spaces.
\end{lemma}

\begin{proof}
Let $V$ be a scheme over $S$ and $y \in \mathcal{Y}_V$. We have to prove
$(\Sch/V)_{fppf} \times_{y, \mathcal{Y}} \mathcal{X}$ is representable
by an algebraic space.

\medskip\noindent
Case I: $V$ is affine and maps into an affine open $\Spec(\Lambda) \subset S$.
Then we can write $V = \lim V_i$ with each $V_i$ affine and of finite
presentation over $\Spec(\Lambda)$, see
Algebra, Lemma \ref{algebra-lemma-ring-colimit-fp}.
Then $y$ comes from an object $y_i$ over $V_i$ for some $i$ by assumption (1).
By assumption (3) the fibre product
$(\Sch/V_i)_{fppf} \times_{y_i, \mathcal{Y}} \mathcal{X}$ is representable
by an algebraic space $Z_i$. Then 
$(\Sch/V)_{fppf} \times_{y, \mathcal{Y}} \mathcal{X}$ is representable
by $Z \times_{V_i} V$.

\medskip\noindent
Case II: $V$ is general. Choose an affine open covering
$V = \bigcup_{i \in I} V_i$ such that each $V_i$ maps into an affine open
of $S$. We first claim
that $\mathcal{Z} = (\Sch/V)_{fppf} \times_{y, \mathcal{Y}} \mathcal{X}$
is a stack in setoids for the Zariski topology. Namely, it is a stack in
groupoids for the Zariski topology by
Stacks, Lemma \ref{stacks-lemma-2-product-stacks-in-groupoids}.
Then suppose that $z$ is an object of $\mathcal{Z}$ over a scheme $T$.
Denote $g : T \to V$ the morphism corresponding to the
projection of $z$ in $(\Sch/V)_{fppf}$. Consider the Zariski sheaf
$\mathit{I} = \mathit{Isom}_{\mathcal{Z}}(z, z)$. By Case I we see that
$\mathit{I}|_{g^{-1}(V_i)} = *$ (the singleton sheaf). Hence
$\mathcal{I} = *$. Thus $\mathcal{Z}$ is fibred in setoids. To finish
the proof we have to show that the Zariski sheaf
$Z : T \mapsto \Ob(\mathcal{Z}_T)/\cong$ is an algebraic space, see
Algebraic Stacks, Lemma
\ref{algebraic-lemma-characterize-representable-by-space}.
There is a map $p : Z \to V$ (transformation of functors) and by Case I
we know that $Z_i = p^{-1}(V_i)$ is an algebraic space. The morphisms
$Z_i \to Z$ are representable by open immersions and
$\coprod Z_i \to Z$ is surjective (in the Zariski topology).
Hence $Z$ is a sheaf for the fppf topology by
Bootstrap, Lemma \ref{bootstrap-lemma-glueing-sheaves}.
Thus Spaces, Lemma \ref{spaces-lemma-glueing-algebraic-spaces}
applies and we conclude that $Z$ is an algebraic space\footnote{
To see that the set theoretic condition of that lemma is satisfied
we argue as follows: First choose the open covering such that
$|I| \leq \text{size}(V)$. Next, choose schemes $U_i$ of size
$\leq \max(\kappa, \text{size}(V))$ and surjective \'etale morphisms
$U_i \to Z_i$; we can do this by assumption (2) and
Sets, Lemma \ref{sets-lemma-bound-size-fibre-product}
(details omitted). Then
Sets, Lemma \ref{sets-lemma-what-is-in-it}
implies that $\coprod U_i$ is an object of $(\Sch/S)_{fppf}$.
Hence $\coprod Z_i$ is an algebraic space by
Spaces, Lemma \ref{spaces-lemma-coproduct-algebraic-spaces}.
}.
\end{proof}

\end{document}
